\section{Mechanism parameters}
\label{app:parameters}

\Cref{tab:parameters} is generated directly from the machine-readable table the
model loads, so it cannot drift from the values the results were produced with.
Each row carries an evidence field and a provenance field in the source file; the
loader refuses to run if a row marked \texttt{sourced} carries no evidence, or if
a row marked \texttt{knob} claims to have been verified.

\begin{table}[h]
\centering
\footnotesize
% Generated by scripts/build_paper_results.py from
% data/reference/mechanism_parameters.csv. Do not edit by hand.
\begin{tabular}{@{}llrrl@{}}
\toprule
Parameter & Kind & Value & Range & Unit \\
\midrule
mainstream persistence & sourced & 0.150 & 0.050--0.300 & correlation \\
mainstream propensity cv & sourced & 0.570 & 0.440--0.800 & coefficient of variation \\
group retention haredi & sourced & 0.867 & 0.735--0.950 & proportion \\
group fertility haredi & sourced & 6.400 & 5.800--7.000 & children per woman \\
group share haredi & sourced & 0.139 & 0.120--0.150 & proportion of population \\
group retention amish & sourced & 0.845 & 0.800--0.920 & proportion \\
group fertility amish & sourced & 6.100 & 5.000--6.500 & children per woman \\
group share amish & sourced & 0.00115 & 0.00110--0.00130 & proportion of population \\
\midrule
mainstream types & knob & 3 & --- & count \\
defector high propensity weight & knob & 0.600 & 0.300--0.900 & proportion \\
development decline per decade & knob & 0.040 & 0.000--0.080 & proportional fall \\
development floor & knob & 0.550 & 0.400--0.750 & multiplier \\
group convergence per generation & knob & 0.100 & 0.000--0.250 & proportional fall \\
\bottomrule
\end{tabular}

\caption{All \MechanismParameterCount{} mechanism parameters. Of these,
\MechanismSourcedCount{} are estimated from published sources and checked against
them, and \MechanismKnobCount{} are scenario knobs with no independent support.
The range column is propagated through the parameter ensemble of
\Cref{sec:results} and the boundary grid of \Cref{sec:boundary}; it is a stated
plausible interval across settings, not a confidence interval from a single
study.}
\label{tab:parameters}
\end{table}

\subsection*{Sources for the estimated rows}

\begin{description}[leftmargin=1.2em,style=nextline,itemsep=0.35em]
\item[mainstream persistence] Parent--child correlation in completed family size.
\citet{murphy1999} reports values mostly 0.15--0.20 in the larger recent British
and American samples reviewed; \citet{murphy2013}, covering 46 populations in 28
countries, reports unweighted regional means from 0.09 to 0.19. The interval
0.05--0.30 is a deliberately conservative cross-setting range rather than one
study's confidence interval. This is combined genetic and cultural continuity,
not a heritability.

\item[mainstream propensity cv] Computed from \CVAllCountries{} pinned country
files of the Cohort Fertility and Education database \citep{cfe2017}; among the
\CVCountries{} with mean completed fertility at or below 2.2 children, the
country-median coefficient of variation is \CVMedian{} and the observed range is
\CVMin{}--\CVMax{}. Independently cross-checked at \CVUnitedStates{} for United
States cohorts 1947--1956 \citep{cdccohort}. See \Cref{sec:cv} for the tail
treatment and the declared inclusion rule.

\item[group retention/fertility/share, haredi] \citet{regev2020} estimate that
13.3\% of people aged 20--64 raised Haredi had left, implying \HarediRetention{}
retention, with documented faction rates from 5.4\% to 26.4\% --- which is what
the wide interval reflects. Fertility \HarediFertility{} is the national
statistical office's own estimate of 6.38 for 2020--2022, which \citet{idi2024}
round to 6.4 and set against a fall from 7.5 in 2003--2005; the same report gives
1.392 million Haredim in 2024, or \HarediShare{} of Israel's population.

\item[group retention/fertility/share, amish] \citet{anderson2025} analyse more
than 50{,}000 households and estimate an age-40+ retention of \AmishRetention{},
which the authors identify as the best available approximation to lifetime
retention, and a total fertility rate of \AmishFertility{} from a 2002--2011
snapshot. Population share \AmishShare{} from \citet{youngcenter2024}, whose
definition covers horse-and-buggy affiliations and excludes Beachy Amish and
Amish Mennonites; the model inherits that boundary.
\end{description}

\subsection*{Why the knobs stay knobs}

No literature check can convert a structural choice or an explicitly future path
into an estimate. The number of mainstream types is a discretisation. The routing
of group leavers has no direct measurement. The rate at which a group's fertility
advantage decays as it stops being unusual has documented instances but no
established general value. And the two environment parameters describe a future
that has not happened.

The last of these is the reason \Cref{sec:boundary} exists. Fitting the
environmental decline to observed fertility would consume exactly the evidence
the mechanism is meant to explain, and setting it by judgement would let it
control a headline it cannot support --- as \Cref{fig:ladder} shows it doing,
at \LadderPressureLoss{} billion against the mechanism's
\LadderMainstreamGain{} billion.
